\chapter*{Smoggy Dues}

\begin{minipage}[t]{\textwidth}
\lettrine{E}rnest Olber stumbled into Willard Jones’s downtown office, his coat dank and face flushed from infusing the choking smog outside. Coughing violently, waving a hand to clear the lingering haze from his nose, the sharp sting of pollution clawing at his throat he reaches Willard’s desk, grabbing the whiskey decanter without waiting for an offer and pours himself a first stiffener.

\begin{figure}[H]
\includegraphics[width=1.0\linewidth]{Illustrations/vign-smog.jpg}
\caption{Olber looking to light up Willard's day}
\end{figure}


Willard typically unruffled, looks up from his notebook, filled as it is with intricate lattices of numbers and half-written equations, raises an eyebrow, and melodes "Rough evening, Ernest?"
\end{minipage}

"Rough is the word," Ernest rasps, knocking back the whiskey in one go, setting down all too dramatically the glass with a unnecessarily jarring clink. "Walking through that smog felt like wading through a cosmic soup. Car headlights stabbing me through the mist, barely visible; it all minded me of how little Baryonic matter there actually is in the universe. Just a thin scattering of stars, photons, and neutrinos compared to the vast, invisible smog of dark matter and that all-encompassing fog of dark energy."

Willard chuckling as his role dictates, pencil tapping now against the edge of his desk ponders whether to bait him. "Baryons are as the primes of the cosmos, scattered and rare, while the dark sector looms vast and undefined?"

Ernest nodding, still pretty breathless. "Exactly. Primes ... visible, important, but nothing compared to the infinite shadow of composites. The universe’s mysteries hide in the spaces between the stars Willard. Just like dark matter and energy dominate the cosmos, the unspoken rules in your lattice is where the mystery lies."

Willard no longer slouching, intriguing himself. "Well, Ernest, it seems the smog outside has clouded your judgment. Tell me more."

Ernest mirroring Willard uncontrollably leans back in his chair, gazing thoughtfully at the night sky. The soft hum of the city a comforting background. "Willard," Ernest began, "you’ve always marveled at the primes. Tell me, do they not ever remind you of stars?"

Willard paused, his pencil hovering midair. "How so?"

Ernest smirking. "Think about it. Primes are scattered across the number line, seemingly random yet governed by deep, immutable laws. Stars, too, are scattered across the cosmo with each a luminous point obeying Oppenhiemer's and Hoyle's laws of stellar nuclear fusion. Their brightness, much like the density of primes, wanes with distance, governed by the inverse square law. But there’s more there, some hidden structures, harmonies that bind them both."

Willard set his notebook aside, intrigued. "Go on, Ernest. Give this one a go."

"Take the power-law relationships of stellar brightness," Ernest said, his voice taking on the tone. "A star’s luminosity, for much of its life, scales with its mass. Larger stars burn brighter, but they burn through their fuel far faster. It’s a fleeting brilliance, much like the sharp spikes of rare primes ... here for a moment, then gone."

Willard chuckling. "And their contributions, like those primes, are singular... and yet the structure of the cosmos as well as that of mathematics rests on these peaks."

Ernest nodding. "Exactly. Now consider this. As light from stars travels to us, its brightness diminishes not just with distance but also with the universe’s expansion. Redshift stretches their wavelengths, thinning the energy of their photons. It’s like adding a decay term to an infinite series. The signal weakens."

Willard still entertaining this flakey allegory, "So, you're saying that interpreting the light of ancient stars is like unraveling the convergence of a series with additional corrections. Each term carries information, but their contributions in turn are diminishing."

"Precisely," Ernest proudly. "And speaking of series, ... have you ever considered how the Riemann zeta function appears in the physics of blackbody radiation?"

Willard now on high BS alert. "In the physics of the cosmos? You’ve managed to merge the zeta function with the cosmic microwave background?"

"Of course!" Ernest’s grin widened. "The zeta function appears naturally when we calculate the number density of photons in a blackbody radiation field. The integration over energy states yields terms like 

\(2\xi(3)\) for photons, or \(\frac{3}{2}\xi(3)\) for neutrinos. It’s all connected, Willard. The lattice of primes and the dance of particles are manifestations of the same underlying harmonies." Willard was silent for a moment. "So," he said finally, "the universe is a lattice of light, governed by rules that echo those of the primes is a cosmos stitched together by zeta-like functions."

\begin{figure}[H]
\includegraphics[width=1.0\linewidth]{Illustrations/vign-blanket.png}
\caption{Ernest at peace with the universe}
\end{figure}

Ernest nodding. "And just as the zeta function connects primes to the infinite, the physics of light attenuation connects us to the stars of the cosmos."

Willard with his pencilset back in motion, "Ernest, perhaps you’ve given me more than just a metaphor."

"And you, Willard," Ernest replying, "have reminded me of how much more elegant the infinite is over the poultry finiteness of the universe. Now, shall we tackle that appendix together for the paper..."

The day's exertions catching up, smog-filled excursions, the whiskey sharpener or two were draining him. Eyes a flutter, shutting as car horns outside fade into the quiet hum of the office. Head now lolling back against the chair, the soft rumble of snoring tremors the room.

Willard all bemused, "Well, old friend," muttering, setting his pencil down, "it seems you’ve found a more immediate comfort." Quietly, he rises from his desk walks to a cabinet, pulling out a soft blanket. The office had cooled considerably despite Ernest’s imposing emanations. Willard knew the night chill would creep in ever so gently, sp he draped the blanket over his friend’s broad shoulders, tucking it around him as best he could. Ernest shifting slightly, murmuring something incomprehensibly satisfying, his face softening as the warmth enveloped him.

"You’re welcome," Willard said under his breath, amused by the faint smile on Ernest’s face. He stepped back and surveyed the scene: Ernest, the cosmologist, out like a light.

The late hour weighed on Willard too as he gathered his papers and notes. The streets outside were quieter now, the smog-laden air muffling the occasional streetcar horn. Willard hesitated, then decided to retreat to the staff common room. 

At this time of the day, an earnest, if rather delicious young woman who always lingered too long in his presence wouldn’t be there now to make him feel awkward. He cast one last glance at his friend, snug under the blanket, whispering "Rest well, Ernest. Tomorrow, the zetas and the smog will still be waiting."

With a soft click, Ernest closed the door behind him, casting a last glance at his old friend, now completely at peace. The streetlights also cast penumbric dances across the room, and for a brief moment, the cosmos seemed a little forgiving.

\begin{figure}[H]
\includegraphics[width=1.0\linewidth]{Illustrations/vign-willardComm.png}
\caption{Willard imbued with a flux of light}
\end{figure}

Clutching Titchmarsh’s \href{https://global.oup.com/academic/product/the-theory-of-the-riemann-zeta-function-9780198533696?cc=gb&lang=en&}{<<The Theory of the Riemann Zeta-Function>>} under his arm, its well-worn spine a testament to countless hours of thunbing, Willard retreated to the staff common room. For now at least, there was no-one to disturb the tranquility, and Willard settled into his usual chair guilt-free.


% >>>>>> ANNEX D (cut-sheet v2): the chemical potential of the universe and the drip-to-equation epilogue — block commented out below, pointer left in text <<<<<<
\textit{(Why photons pay no chemical rent, and how a whiskey drip becomes an equation, is worked in Annexe D.)}

% \section*{Postscript: Chemical potential of the Universe}
% 
% 
% Calculating \(\mu\) provides insight into energy per particle across the various particle species. Here we focus on relativistic particles. For a photon gas in thermal equilibrium, the internal energy \( U \) can be expressed in terms of temperature \( T \), volume \( V \), and photon number \( N_\gamma \). To see this consider again the Stefan-Boltzmann law ~\ref{eqSB}, from which the energy density \( \rho_\gamma \) reads:
% \[
% \rho_\gamma = \mathcal{A} T^4,
% \]
% \noindent
% where \( \mathcal{A} = \frac{4 \sigma}{c} \). The total internal energy \( U \) in a volume \( V \) is:
% \begin{equation}
% U = \rho_\gamma V = \mathcal{A} T^4 V.\label{eq:tee4}
% \end{equation}
% \noindent
% The photon number density \( n_\gamma \) scales as \( T^3 \):
% \begin{equation}
% n_\gamma = b T^3, \label{numDense}
% \end{equation}
% \noindent
% where \( b = \frac{16 \pi \zeta(3) k_B^3}{c^3 h^3} \). For total photon number \(
% N_\gamma = n_\gamma V = b T^3 V.
% \)
% we can solve for \( T \):
% \[
% T = \left( \frac{N_\gamma}{b V} \right)^{1/3}.
% \]
% \noindent
% Substituting this expression into(\ref{eq:tee4}) yields:
% \[
% U = \mathcal{A} b^{-4/3} N_\gamma^{4/3} V^{-1/3}.
% \]
% \noindent
% The chemical potential for a photon thus reads:
% \begin{align*}
% \mu_\gamma &= \left( \frac{\partial U}{\partial N_\gamma} \right)_{S, V} \\
% &= \frac{\partial}{\partial N_\gamma} \left( \mathcal{A} b^{-4/3} N_\gamma^{4/3} V^{-1/3} \right) \\
% &= \frac{4}{3} \mathcal{A} b^{-4/3} N_\gamma^{1/3} V^{-1/3},\\
% & \approx 2.0 \times 10^{-17} \, \text{J}.
% \end{align*}
% 
%%%%%%%%%%%%%%%%%%%%%%%%
% which includes within \(b\) the Riemann zeta function, \(\zeta(s)\), for \(s=3\). To see where this arises in (\ref{numDense}) consider integrating the distribution functions for both photon and neutrino gases across all energy states, reflecting the equilibrium behavior of bosons and fermions. The number density \( n \) of blackbody radiation field and neutrino gases is obtained by integrating the Bose-Einstein distribution (resp. Fermi-Dirac distribution) over all energy states:
% \begin{align*}
% n_\gamma &= \int_0^\infty \frac{g(\epsilon)}{e^{\epsilon / k_B T} - 1} \, d\epsilon, \quad
% n_\nu = \int_0^\infty \frac{g(p)}{e^{p c / k_B T} + 1} \, dp,
% \end{align*}
% \noindent
% where \( g(\epsilon) = \frac{8 \pi \epsilon^2}{(hc)^3} \), \( g(p) = \frac{4 \pi p^2}{h^3} \). Changing variables with \( x = \frac{\epsilon}{k_B T} \) and \( x = \frac{p c}{k_B T} \), respectively, we find:
% 
% \begin{align}
% n_\gamma &= \frac{8 \pi (k_B T)^3}{(hc)^3} \int_0^\infty \frac{x^2}{e^x - 1} \, dx,\label{numbPho}\\ 
% n_\nu &= \frac{4 \pi (k_B T)^3}{h^3 c^3} \int_0^\infty \frac{x^2}{e^x + 1} \, dx.\notag
% \end{align}
% The two integrals over the dimensionless variable \( x \) yield \(2\zeta(3) \) and \(\frac{3}{2}\zeta(3)\)
% \[
% \int_0^\infty \frac{x^2}{e^x - 1} \, dx = 2 \zeta(3) \text{,}\quad\int_0^\infty \frac{x^2}{e^x + 1} \, dx = \frac{3}{2} \zeta(3),
% \]
% giving the final results:
% \begin{align}
%     n_\gamma &= \frac{16 \pi \zeta(3) (k_B T)^3}{(hc)^3},  \label{eq:numbP} \quad
%     n_\nu  = \frac{6 \pi \zeta(3) (k_B T)^3}{(hc)^3}. \notag
% \end{align}
% To source the \(\zeta(3)\) in (\ref{eq:numbP}) expand the integral (\ref{numbPho}) as a geometric series, 
% \begin{align*}
%     \int_0^\infty \frac{x^2}{e^x - 1} \, dx &= \int_0^\infty x^2 \sum_{n=1}^\infty e^{-nx} \, dx\\ &= \sum_{n=1}^\infty \int_0^\infty x^2 e^{-nx} \, dx.
% \end{align*}
% \noindent
% Noting the integral form is akin to the gamma function, \[
% \Gamma(z) = \int_0^\infty t^{z-1} e^{-t} \, dt, \]
% for \( z = 3 \) with \(\Gamma(3) = 2!\), we set \( t = nx \), so \( x = \frac{t}{n} \) and \( dx = \frac{dt}{n} \). 
% 
% \noindent
% We have then that
% \begin{align*}
%     \int_0^\infty x^2 e^{-nx} \, dx &= \frac{1}{n^3} \int_0^\infty t^2 e^{-t} \, dt\\\
%     &= \frac{1}{n^3} \Gamma(3) = \frac{1}{n^3} \cdot 2 = \frac{2}{n^3}.\\
%     \int_0^\infty \frac{x^2}{e^x - 1} \, dx &= 2 \sum_{n=1}^\infty \frac{1}{n^3} = 2 \zeta(3).
% \end{align*}
% 
%%%%%%%
% \newpage
% \section*{Epilogue: From Drip to Equation}
% 
% \begin{quote}
% \emph{“The proper place for play—the sweet spot of life—is the dancing edge,  
% the curious middle ground between too much order and too much chaos.”}  
% \hfill — Alan Watts (1966)
% \end{quote}
% 
% The call began, as these things often did, with a leak. Ernest sitting there in half-thought as a single cold drop of water eeks out a path through a cracked light fitting striking the crown of his head. Then a a slow, perceptibly deliberate drip. Creations more impressive than any previously achieved on his porcelain throne. Conjured and polished. That graceful fall, a parable of confluence over effluence.
% 
% Not a full blown leak exactly. A drop or two,.. three... Not ok. But there was something oddly faithful about the replication, the water remembering its path, tracing a martingale of descent.
% 
% Later that day, he called Willard, telling him it wasn’t the drop that mattered but its return path:  how water, finding the weakest line in the ceiling, had discovered the only route it could.
% 
% ``It wasn’t recurrence,'' he said. ``It was selection by failure.  The leak wasn’t remembering its path, rather it was obeying it.''
% 
% Willard chuckled, used to Ernest’s inversions.
% 
% But Ernest pressed on, the sound of kettle steam hissing somewhere behind him.
% 
% ``That’s what confluence is, Ed. We think we choose, we branch, we bifurcate in Everett’s forking garden but what if most of those branches collapse?  
% What if persistence is just the world finding its low resistance groove?  
% Teleology without foresight.  Purpose as the path of least entropy.''
% 
% He then spoke of his life as if he were a model.  Not just Prada, a Platonic build comprised of periodic aspirations: those jobs, relationships, papers all that occasionally flared as strange attractors before inexplicably waning in their want.  Each time he thought he had actively changed direction -all that built-in inertia and memory bringing him back into a regulated orbit- suddenly he would be off again.  Yet perhaps, he bemused, that need for change wasn’t a signal of failure.  Perhaps that persistence to change, the adaptive re-weighting of purpose under noise was his build
% 
% ``Maybe,'' Ernest said, half to himself, ``we never escape our Markovian chains.  
% We just shift the probabilities until they start to look like pure will.''
% 
% Willard smiled on the other end.
% 
% ``And this is what you’re coding?''
% 
% Ernest nodded though the line was silent.  
% He’d already opened the notebook---\href{https://colab.research.google.com/drive/1sKF2IJpjPXPKs1o3Ap06QmU7RzpcHo7K?usp=sharing}{Teleological Markov Model}.  
% The leak had the quality of an inequality:  
% a retarded force of memory, with an advanced pull of anticipation, and a noise term large enough to make both inexpplicableto all but the most curious.  
% He would call this paper \emph{On Confluence Without Collapse}.
% 
% ``The universe,'' he said affectedly, ``doesn’t need foresight.  
% It just needs persistence enough to look like it is in possession of purpose.''
% 
% \medskip
% The call ended, the room returning to its old rhythm.  An oh so deliberate drip making faithful, fateful return.
% 
% He smiled. On the screen, lines of Python flickered like synapses:  
% $\lambda(t)$, $\tau(t)$, $C_i(t)$.  
% It was all there: a model of persistence, inertia, and noise; a conjecture that even stochastic wanderers can appear to find some kind of purpose.
% 
% \bigskip
% 
% He’d written it once before, in a now discarded draft:  
% ``Purpose is to the memory of future states, as habit is the inertia of past ones.''
% 
% \bigskip
% Now, hearing the drip, he actually began believing it.  Every recurrence a negotiation between what is trying to endure and what is yielding.
% 
% ``Maybe,'' Ernest said softly, ``the universe doesn’t aim; rather it remembers.  
% And we only call that memory ‘direction.’''
% 
% Outside, the light flickered in from a late sky, and on the screen the simulation compiled.  He would call the notebook \emph{Enhanced Markov Confluence and Time-Symmetric Potentials}.  It would begin with drip rhythm and end with the thought that coherence, might be nothing more than a well-timed release.
% 
% Willard recalled, ``As Alan Watts\cite{Watts1966} observed, the vitality of both games and life lies in their “curious middle ground”—the dancing edge between too much order and too much chaos.  ’''
% 
% Ernest's hybrid model was trying to formalize this balance: between retarding force anchors of stability and advancing force inviting surprise.  When tuned in proportion, they sustain. Keeping up a playfulness of coherence without determinism, purpose without prescription.
% 
% \vspace{0.6em}
% \noindent
% The next morning, the note written in a Jupyter cell became the opening line of this study:
% 
% \begin{quote}
% \emph{Systems appear to want what they are structurally tuned to predict.}
% \end{quote}
% 
% % Requires: \usepackage{tcolorbox}
% \begin{tcolorbox}[colback=gray!5,colframe=gray!40,title={Jupyter cell: Ernest’s note to self}]
% {\small\ttfamily
% \# Notebook: On Apparent Purpose\\[4pt]
% \#\\
% \# Observation:\\
% \# Systems appear to *want* what they are structurally tuned to *predict*.\\
% \# Teleology is the residue of temporal correlation ---\\
% \# a forward echo mistaken for desire.\\[6pt]
% \#\\
% \# When $\textbackslash lambda(t)$ rises, the present leans into its own forecast.\\
% \# The drip of expectation becomes a groove,\\
% \# and coherence emerges not from choosing,\\
% \# but from remembering where the next drop is likely to fall.
% }
% \end{tcolorbox}
% 
% From that the equations followed.  Just as droplets trace paths of least resistance, so too do stochastic agents carve recurrent trajectories through noise: not by foresight, but by a feedback between memory and anticipation.  
% Persistence arises not as opposition to stochasticity but as by submissive means.
% 
% \vspace{0.6em}
% \noindent
% In this sense, the model is less about purpose than about the appearance of purpose:  
% how forward-weighted correlations, when coupled with bounded recollection, 
% give rise to behavior that looks intentional.  Teleology is not to be a metaphysical claim but an emergent geometry of feedback:%
% a system learning to repeat what once cohered.
% 
% \vspace{0.6em}
% \noindent
% As Ernest writes in his notebook:
% 
% \begin{quote}
% \emph{The world does not conspire toward purpose, yet recurrences make it seem so.  
% We name that appearance ‘meaning,’ though it is only the confluence of what remembers, endures  and what expects.}
% \end{quote}
% 
% \vspace{0.6em}
% \noindent
% Thus a convergence on a single insight:  
% that coherence may be the most faithful expression of purpose available to a stochastic world: a quiet geometry of purposeful wandering.\footnote{%
% See accompanying notebook: \href{https://colab.research.google.com/drive/1MA6otuq2OctiqEkDxIH21gV9eOUhYhv_?usp=sharing#scrollTo=uvGmJnfw69Lx}{\texttt{Enhanced Markov Confluence and Time-Symmetric Potentials}}.  
% The simulations replicate the dynamics discussed in this paper, linking the narrative reflection to empirical visualization.
% }
% 
% \vspace{0.8em}
% \noindent
% Colleagues would likely refer to it as \emph{“An Ernest Armitage shanks epiphany.”} Perhaps it was all shank and no drive, but he enjoyed retracing its conjure.
% 
% 
% % Requires in preamble: \usepackage{tcolorbox}
% \begin{tcolorbox}[colback=gray!5,colframe=gray!40,title={Ernest’s Overleaf preamble}]
% {\small\ttfamily
% \textbackslash title\{\textbackslash textbf\{Markovian Confluence and the Geometry of Purpose\}\textbackslash\textbackslash[4pt]\\
% What began as a ceiling leak became a model of purpose: a feedback loop between recollection, anticipation, and chance\}\\[6pt]
% 
% \textbackslash author\{Ernest Olber\textbackslash\textbackslash\\
% \textbackslash Department of Applied Epistemics, Merriweather Institute\textbackslash\textbackslash\\
% \textbackslash \textbackslash href\{mailto:ernest.olber@institute.edu\}\{ernest.olber@institute.edu\}\}\\[3pt]
% 
% \textbackslash date\{\textbackslash today\}\\[6pt]
% 
% \textbackslash maketitle\\[8pt]
% 
% \textbackslash vspace\{1em\}\\
% \textbackslash paragraph*\{Author’s Note.\}\\
% This model originated from a moment of unplanned reflection---\%\\
% a single drop of water falling through a ceiling onto the author’s head.
% }
% \end{tcolorbox}
% 
% \textbf{Willard} grinning. ``So this is just the owl and cheese all over again, isn’t it\cite{FriedmanForster2005}
% Give the mouse a prize and it learns the maze; give it an owl and it freezes.  
% Your $\lambda(t)$ is the cheese term: turn it up and everyone moves faster.''
% 
% \textbf{Ernest} shook his head. ``Not quite. That mouse experiment plays with motivation: affective framing, approach versus avoidance.  
% What I’m modeling isn’t emotion, it’s time structure.  
% $\lambda(t)$ doesn’t sweeten the goal; it changes how much of the future enters the present equation.  
% Cheese excites the limbic system while $\lambda(t)$ modulates temporal coupling.  
% One’s psychological, the other mechanical.''
% 
% He tapped a finger on the terminal where the trajectories pulsed across the screen.
% 
% ``My agents don’t hope, they integrate.  
% If they seem purposeful, it’s because anticipation reduces entropy, not because they’ve tasted reward.  That’s the difference between cognition and confluence, Ed:
% the system behaves as if it wants, even though all it’s really doing is weighting tomorrow a little more than yesterday.''
% 
% Ernest turned back to his screen.
% \smallskip
% ``Here, Ed .... this is the difference.  
% No cheese, no owl, no dopamine ... just weighted foresight.''
% 
% \smallskip
% 
% The screen drew its familiar dance: a red trace bending, hesitating, then following a dashed black trail that drifted lazily, cloud-like.
% 
% \smallskip
% 
% ``You see?'' Ernest said. ``The agent isn’t chasing a prize.  
% It’s just tuning the future into the present.  
% As $\lambda$ rises, anticipation dominates inertia. That is coherence from coupling, not craving.''
% 
% \smallskip
% 
% Willard watched the looping paths, the convergence without command.  
% ``It still looks,'' he said, ``like it wants the cheese.''
% 
% Ernest smiled. ``Only because you do.'' He found consolation in the \href{https://leelemacjames.github.io/MarkovConfluence/}
% {firmaments} of his console.
% 
% \begin{figure}[H]
% \includegraphics[width=1.0\linewidth]{Illustrations/console.png}
% \caption{Willard's Firmament consolations}
% \end{figure}
% 
% \href{https://www.researchgate.net/publication/396515722_Markov_Confluence_and_Time-Symmetric_Potentials_Persistence_not_proliferation}{The Tug of War} in which every agent lives in tension: memory whispers "stay close to where you've been" while teleology beckons "move toward where you're going." The dial\(\lambda(t)\) arbitrates this quarrel at dawn, the past wins; by dusk, the future pulls harder. What emerges isn't paralysis but purposeful wandering: trajectories that neither forget their origins nor ignore their destinations. Persistence, it turns out, is not about choosing sides, rather it's about knowing when to listen to each voice.
% 
% \begin{tcolorbox}[colback=gray!5,colframe=gray!40,title={Jupyter cell: Weighted foresight simulation}]
% {\footnotesize\ttfamily
% import numpy as np\\
% import matplotlib.pyplot as plt\\[4pt]
% T, dt = 200, 0.1\\
% x = np.zeros((T, 2))~~~~~\# agent position\\
% g = np.zeros((T, 2))~~~~~\# moving goal\\
% v = np.zeros((T, 2))~~~~~\# velocity\\
% lam = np.linspace(0, 1, T)~~~~~\# teleological gain schedule\\
% a\_ret, a\_adv, gamma = 0.25, 0.12, 0.6\\[4pt]
% for t in range(1, T):\\
% ~~~~g[t] = g[t-1] + 0.1*np.random.randn(2)~~~~~\# Brownian goal drift\\
% ~~~~F\_ret = -a\_ret * (x[t-1] - np.mean(x[:t], axis=0))~~~~~\# memory pull\\
% ~~~~F\_adv = -lam[t] * a\_adv * (x[t-1] - g[t])~~~~~\# anticipatory pull\\
% ~~~~v[t] = gamma*v[t-1] + dt*(F\_ret + F\_adv) + 0.1*np.random.randn(2)\\
% ~~~~x[t] = x[t-1] + v[t]\\[4pt]
% plt.plot(x[:,0], x[:,1], 'r-', label='agent')\\
% plt.plot(g[:,0], g[:,1], 'k--', label='goal')\\
% plt.legend(); plt.axis('equal'); plt.show()
% }
% \end{tcolorbox}
% 
% 
% 
% \newpage
% 
%%%%%
% (Willard Jones' Flux chapter moved to 8-15willardFlux.tex — cut-sheet v2)
% >>>>>> end of annexed block <<<<<<
